\section{A NoSQL Projection}
\label{sec:annex-nosql}

The \texttt{sfs-nosql} crate is an intentionally secondary experiment: if a
unit need not be a pathname, can the fragment-causal engine support a document
interface without becoming a second storage system?  It represents a record as
a store name, a 128-bit primary key, and a sorted map from property names to
typed values.  The value codec supports null, Boolean, signed 64-bit integer,
IEEE-754 double, UTF-8 string, and byte-string values.  Records are ordinary
\sfs{} units under an internal keyspace, so they reuse engine publication,
history, encryption, and synchronization records.

The projection implements put, get, delete, primary-key listing, and batched
\texttt{put\_many}.  The batch executes all record and index updates under one
engine transaction and one header publication.  Secondary index units map
\((\mathit{store},\mathit{property},\mathit{value})\) to sorted primary-key
sets, supporting equality and inclusive same-type range queries.  Large opaque
values can use \texttt{put\_unindexed} and remain addressable by primary key
but are intentionally absent from secondary queries.

The more informative experiment is property-granular merge.  A patchable record
uses a 4-KiB directory slot followed by one aligned 4-KiB slot per property.
Changing an existing property rewrites its slot and the affected secondary
index entries in one transaction, leaving the causal dots of other property
fragments unchanged.  The artifact scenario starts two replicas from one
record, changes property \texttt{alpha} on one and \texttt{beta} on the other,
then imports the changed fragment: both values appear in one strain.  Repeating
the scenario with whole-record writes produces concurrent strains.  This is a
direct demonstration that application layout can expose semantic independence
to the same fragment merge rule used for files.

The result is not a general NoSQL claim.  Records are flat rather than nested,
there is no query planner, distributed transaction protocol, or database
benchmark, and indexed values must fit the internal trie-key bound.  The current
slotted scheme is especially narrow: the directory and every value must each
fit 4 KiB, and the complete image must remain below the 16-KiB geometry
transition.  With one directory slot, that permits at most two property slots
in the present implementation.  Adding a property rewrites the directory and
requires a full patchable put.

Finally, the merge scenario drives record and fragment import directly.  It does
not run the full \texttt{SyncEngine} or establish convergence of secondary
indexes after two remote property patches.  Thus the annex supports one
resolution claim: a fragment-causal backend can be projected into
property-granular merge by layout.  Distributed index repair and a production
document model remain future work.
